% ============================================================
%  SLIDES EXPÉRIMENTALES — à insérer dans tipe_cogeneration.tex
%  REMPLACE les deux slides "Perspectives expérimentales (1/2)" et
%  "Perspectives expérimentales (2/2)" (lignes 1308–1409)
%
%  Moteur Stirling Jeulin 204081 — ExAO Foxy — Jérémie Prouteau
%
%  Nécessite dans le préambule :
%    \usepackage{graphicx}   (déjà présent)
%    \usepackage{pgfplots}   (déjà présent)
%    \pgfplotsset{compat=1.18}
%
%  Figures à placer dans fig/ :
%    fig/stirling_PV.pdf        ← diagramme P-V mesuré (script Python)
%    fig/stirling_labo.jpg      ← photo du montage (vos photos)
%    fig/stirling_eta_comp.pdf  ← graphe η_exp vs η_Carnot vs η_modèle
%                                  (optionnel — placeholder sinon)
% ============================================================

%% ── Expérience Stirling : dispositif ────────────────────────
\begin{frame}{Validation exp\'erimentale --- Moteur Stirling de laboratoire}
  \framesubtitle{Moteur Jeulin 204081 (3B Scientific) --- ExAO Foxy}

  \begin{columns}[T]
    \column{0.52\textwidth}
      \begin{block}{Dispositif exp\'erimental}
        \small
        \begin{itemize}\itemsep3pt
          \item Moteur Stirling alpha transparent
                (cylindre en verre thermorésistant)
          \item Source chaude : r\'echaud \`a alcool
          \item Capteur PV Stirling Jeulin +
                interface ExAO Foxy~: acquisition
                en temps r\'eel de $P(t)$ et $V(t)$
          \item Thermocouple type~K pour $T_{\text{chaud}}$
                et $T_{\text{froid}}$
        \end{itemize}
      \end{block}

      \vspace{0.4em}
      \begin{block}{Caract\'eristiques du moteur}
        \scriptsize
        \renewcommand{\arraystretch}{1.15}
        \begin{tabular}{@{}ll@{}}
          \toprule
          Param\`etre & Valeur \\
          \midrule
          $V_{\min}$ / $V_{\max}$ & 32 / 44\,cm$^3$ \\
          $\rv = V_{\max}/V_{\min}$ & \textbf{1{,}375} \\
          Puissance nominale & 1{,}5\,W \\
          Vitesse au ralenti & $\approx$\,1\,000\,tr/min \\
          \bottomrule
        \end{tabular}
      \end{block}

    \column{0.44\textwidth}
      \begin{center}
        % ── Remplacer par votre photo ──────────────────────
        \includegraphics[width=0.95\linewidth]{fig/stirling_labo.jpg}
        % \fbox{\rule{0pt}{3.2cm}\rule{3.8cm}{0pt}}
        \\[3pt]
        \small Montage ExAO --- r\'ef.\ 204081
      \end{center}
  \end{columns}

\end{frame}

%% ── Diagramme P-V mesuré ─────────────────────────────────────
\begin{frame}{Diagramme P-V exp\'erimental du cycle de Stirling}
  \framesubtitle{342 points acquis via ExAO Foxy ---
    $T_{\text{démarrage}} = 124\,\text{\textdegree C}$,
    $T_{\text{max}} = 228\,\text{\textdegree C}$}

  \begin{columns}[T]
    \column{0.58\textwidth}
      \begin{center}
        % ── Insérer la figure générée par le script Python ─
        \includegraphics[width=\linewidth]{fig/stirling_PV.pdf}
        % \fbox{\rule{0pt}{4.5cm}\rule{5.8cm}{0pt}}
        \\[-2pt]
        \tiny Diagramme de Clapeyron mesuré (P en mbar, V en cm$^3$)
      \end{center}

    \column{0.38\textwidth}
      \begin{block}{Lecture du diagramme}
        \small
        \textbf{Branche haute} (rouge) :\\
        d\'etente \`a $T_{\max} = 228\,$\textdegree C.\\[4pt]
        \textbf{Branche basse} (bleu) :\\
        compression \`a $T_{\min}$.\\[4pt]
        \textbf{Aire du cycle} :\\
        \[
          \Wnet = \oint P\,dV
        \]
      \end{block}

      \vspace{0.3em}
      \begin{alertblock}{Valeurs mesur\'ees}
        \small
        $\rv^{\text{exp}} = 1{,}368$\\[2pt]
        $\Delta P = P_{\text{haut}} - P_{\text{bas}}$\\
        $\phantom{\Delta P} \approx 246\,$mbar\\[2pt]
        $\Wnet \approx \mathbf{323\,\text{mJ}}$
      \end{alertblock}

      \vspace{0.3em}
      \begin{exampleblock}{}
        \small
        $\etaC = 1 - \dfrac{397}{501} \approx \mathbf{20{,}8\,\%}$\\[2pt]
        {\footnotesize ($T_{\min}$ = temp.\,d\'emarrage)}
      \end{exampleblock}
  \end{columns}

\end{frame}

%% ── Comparaison modèle / expérience ─────────────────────────
\begin{frame}{Confrontation mod\`ele -- mesure}
  \framesubtitle{Identification de l'efficacit\'e du r\'eg\'en\'erateur $\varepsilon$}

  \begin{columns}[T]
    \column{0.50\textwidth}
      \begin{block}{Mod\`ele CoolProp ($\rv = 1{,}368$)}
        \small
        Avec $T_{\min} = 397$\,K, $T_{\max} = 501$\,K,
        $\rv = 1{,}368$ et $\varepsilon$ identifi\'e :
        \[
          Q_{\text{in,eff}} = Q_{34}
            + (1-\varepsilon)\,Q_{\text{23,r\'eg}}
        \]
        \[
          \eta_{\text{mod}} = \frac{\Wnet}{Q_{\text{in,eff}}}
          < \etaC
        \]
      \end{block}

      \vspace{0.4em}
      \begin{block}{Bilan num\'erique}
        \scriptsize
        \renewcommand{\arraystretch}{1.2}
        \begin{tabular}{@{}lc@{}}
          \toprule
          Grandeur & Valeur \\
          \midrule
          $\etaC$ & 20{,}8\,\% \\
          $\eta_{\text{mod}}$ ($\varepsilon = 0{,}85$) & $\approx$\,12\,\% \\
          $\Wnet$ (ExAO) & 323\,mJ/cycle \\
          $\rv^{\text{exp}}$ & 1{,}368 \\
          \bottomrule
        \end{tabular}
      \end{block}

    \column{0.46\textwidth}
      \begin{center}
        % ── Graphe η_exp vs η_Carnot vs η_modèle (optionnel) ─
        \includegraphics[width=\linewidth]{fig/stirling_eta_comp.pdf}
        % \fbox{\rule{0pt}{3.8cm}\rule{4.4cm}{0pt}}
        \\[-2pt]
        \tiny $\eta$ en fonction de $T_{\max}$ ---
              mod\`ele CoolProp vs Carnot
      \end{center}

      \vspace{0.3em}
      \begin{alertblock}{Conclusion exp\'erimentale}
        \small
        Le faible $\rv = 1{,}37$ est le facteur
        \textbf{limitant principal} :
        $\Wnet \propto \ln \rv$ est petit.\\[3pt]
        L'\'ecart $\etaC - \eta_{\text{mod}}$ quantifie
        les irr\'eversibilit\'es r\'eelles
        (r\'eg\'en\'eration imparfaite, frottements).
      \end{alertblock}
  \end{columns}

\end{frame}

% ============================================================
%  FIN DES SLIDES EXPÉRIMENTALES
%  La suite (appendix + annexes) reste inchangée
% ============================================================
